% ══════════════════════════════════════════════════════════════════════════════
%  ADVANCED SPACE — the above-level book (from 2027)
%  Faculty decision, Nov 2026: every student completes the at-level lesson book;
%  students who excel move on to this book. Its topics go deeper into the ideas
%  of the at-level lessons rather than adding more of the same work.
%
%  Topics 1 and 2 are written here around the master's Year 8 extension boxes
%  from Lessons 1, 3, 4 and 6 (%%ABOVE:n%%, n counting them in the order they
%  appear in Lessons 1-7: 1-2 Lesson 1, 3 Lesson 3, 4 Lesson 4, 5 Lesson 6,
%  6 Lesson 7). Topics 3-5 are the master's Lessons 8-10, whole (%%LESSON:n%%);
%  Lesson 7's extension box closes Topic 5. Built by make_a5.py.
% ══════════════════════════════════════════════════════════════════════════════

% ══════════════════════════════════════════════════════════════════════════════
%  TOPIC 1 — THE MOON UP CLOSE                         (deepens Lessons 1, 3, 4)
% ══════════════════════════════════════════════════════════════════════════════
\setcounter{section}{0}
\lessonopener{1}{The Moon Up Close}{../art/l3-ink.pdf}{1 \textbullet\ The Moon Up Close}\label{sec:1}\pdfbookmark[1]{Topic 1}{lesson.1}\hypertarget{lesson.1.1}{}
\vc{Deepens Lessons 1, 3 and 4}{Where the Moon came from; why eclipses are rare; how the Moon's orbit is growing}

\begin{keyvocab}
  \vocabrow{Giant Impact Hypothesis}
  \vocabrow{Isotope}
  \vocabrow{Node}
  \vocabrow{Saros cycle}
  \vocabrow{Angular size}
  \vocabrow{Annular eclipse}
  \vocabrow{Recession}
\end{keyvocab}

\subsection{Where the Moon came from}\label{sub:a1.1}

Every rock has a chemical fingerprint. Elements come in slightly different versions, called \vocab{isotope}s, and the mix of oxygen isotopes in a rock depends on where in the early solar system its material formed. That fingerprint is the key evidence in the reading below.

%%ABOVE:1%%

\begin{youdobox}
\question{2}{Describe the Giant Impact Hypothesis in two sentences.}
\answerbox{1.6cm}

\question{2}{Explain how each of the two pieces of evidence supports it.}
\answerbox{1.6cm}

\question{2}{An earlier idea was that the Moon was a passing body captured by Earth's gravity. Which piece of evidence counts most strongly against this? Explain.}
\answerbox{1.6cm}
\end{youdobox}

\subsection{Why eclipses are rare}\label{sub:a1.2}

%%ABOVE:3%%

How big something looks in the sky is its \vocab{angular size}. It depends on two things: its real size and how far away it is. Two objects look the same size when their diameter divided by their distance is the same.

\begin{idobox}{Do the Sun and Moon really look the same size?}
Sun: diameter 1\,392\,000\,km, distance 150\,000\,000\,km.\quad $1\,392\,000 \div 150\,000\,000 = 0.0093$\\
Moon: diameter 3474\,km, distance 384\,400\,km.\quad $3474 \div 384\,400 = 0.0090$

\smallskip
The two numbers are almost equal, so the Moon only just covers the Sun. When the Moon is a little further away than average, it does not quite cover it, and a ring of Sun shows around it: an \vocab{annular eclipse}.
\end{idobox}

\begin{wedobox}{Eclipses in the far future}
The Moon moves away from Earth by about 3.8\,cm a year (its \vocab{recession}). In 600 million years, how much further away will it be, in kilometres? (1\,km $=$ 100\,000\,cm.) Calculate its new diameter-to-distance number, and decide whether total eclipses of the Sun will still happen.
\workbox{5.0cm}
\end{wedobox}

\subsection{A growing orbit}\label{sub:a1.3}

%%ABOVE:4%%

\begin{youdobox}[Challenge]
\question{4}{Use what you have learnt in this topic to explain why the people living in 600 million years will never see a total eclipse of the Sun, and why the people living 600 million years ago saw them more often than we do. Include numbers.}
\workbox{5.0cm}
\end{youdobox}

% ══════════════════════════════════════════════════════════════════════════════
%  TOPIC 2 — EARTH'S SHIELD AND CROWDED ORBITS             (deepens Lessons 1, 6)
% ══════════════════════════════════════════════════════════════════════════════
\setcounter{section}{1}
\lessonopener{2}{Earth's Shield and Crowded Orbits}{../art/l6-ink.pdf}{2 \textbullet\ Earth's Shield and Crowded Orbits}\label{sec:2}\pdfbookmark[1]{Topic 2}{lesson.2}\hypertarget{lesson.2.1}{}
\vc{Deepens Lessons 1 and 6}{Earth's magnetic field and space weather; the energy of orbital debris; who looks after near-Earth space}

\begin{keyvocab}
  \vocabrow{Dynamo effect}
  \vocabrow{Magnetosphere}
  \vocabrow{Solar wind}
  \vocabrow{Space weather}
  \vocabrow{Kinetic energy}
  \vocabrow{Kessler syndrome}
  \vocabrow{Tragedy of the commons}
\end{keyvocab}

\subsection{Earth's magnetic shield}\label{sub:a2.1}

%%ABOVE:2%%

\begin{youdobox}
\question{2}{Explain where Earth's magnetic field comes from.}
\answerbox{1.6cm}

\question{2}{Why are auroras seen mostly near the poles?}
\answerbox{1.6cm}

\question{2}{Name two kinds of technology that space weather can damage, and why.}
\answerbox{1.6cm}
\end{youdobox}

\subsection{Why a fleck of paint is dangerous}\label{sub:a2.2}

Satellites in low orbit travel at about 7.5\,km every second. At that speed, even tiny pieces of debris carry enormous \vocab{kinetic energy}, the energy of movement. A fleck of paint once chipped a window of a Space Shuttle. In 1978 the NASA scientist Donald Kessler warned that if enough debris builds up, one collision would make fragments that cause more collisions, which make more fragments, until some orbits became unusable. This runaway chain is called the \vocab{Kessler syndrome}.

\begin{idobox}{The energy of a fragment}
\[ \text{kinetic energy (J)} = \tfrac{1}{2} \times \text{mass (kg)} \times \text{speed (m/s)}^2 \]
A 1.4\,g fragment of aluminium (0.0014\,kg) at 7500\,m/s:\\
$\tfrac{1}{2} \times 0.0014 \times 7500^2 = 0.0007 \times 56\,250\,000 \approx \mathbf{39\,000\ J}$\\
That is about the kinetic energy of a one-tonne car at 32\,km/h, all delivered to a spot smaller than a fingernail.
\end{idobox}

\begin{wedobox}{A lost bolt}
A 10\,g bolt (0.010\,kg) drifts away from a spacewalk at orbital speed, 7500\,m/s. Calculate its kinetic energy. How many times more is that than the aluminium fragment?
\workbox{4.5cm}
\end{wedobox}

Removing debris is expensive, and whoever pays for it helps every other satellite owner as much as themselves. Economists call this the \vocab{tragedy of the commons}: when something is shared by everyone, no single owner has a reason to pay to look after it, so it is overused and damaged.

%%ABOVE:5%%

\begin{youdobox}[Challenge]
\question{4}{Explain why the Kessler syndrome could make a crowded orbit unusable even if no one launched anything new. Then propose one rule that every country launching satellites should follow, and explain why it would be hard to get every country to agree.}
\workbox{5.0cm}
\end{youdobox}

% ══════════════════════════════════════════════════════════════════════════════
%  TOPICS 3-5 — THE MASTER'S LESSONS 8-10, WHOLE
% ══════════════════════════════════════════════════════════════════════════════
%%LESSON:8%%

%%LESSON:9%%

%%LESSON:10%%

\subsection{Knowledge that lasts}\label{sub:a5.last}

The astronomy in Lesson 7 was recorded and passed on without writing. How reliable can knowledge passed on that way be over thousands of years? The research below suggests: very.

%%ABOVE:6%%

\begin{youdobox}[Challenge]
\question{4}{Scientists test a claim by checking it against independent evidence. Choose one example from the reading and explain how scientists could check that the oral record was accurate. Then explain why recording knowledge in several ways at once (song, story, ceremony, art) makes it harder to lose or change.}
\workbox{5.0cm}
\end{youdobox}
